\section{Methodology}
\label{sec:method}

We introduce \textbf{Endosymbiotic Holographic Parameter Binding} (EHPB), an interdisciplinary ensembling paradigm that rejects the linear weight-averaging assumption. Rather than directly summing parameters in their native coordinate space, we treat each layer's weight space as a holographic associative memory. In this section, we present the formal mathematical framework of EHPB, detailing key generation, holographic superposition, dynamic routing, parameter transcription, and provide a rigorous proof of its reconstruction accuracy.

\subsection{Mathematical Framework}

Let $W_{\text{base}}^{(l)} \in \mathbb{R}^{R \times C}$ denote the pre-trained, shared base model weight matrix for layer $l$ of a deep neural network, where $R$ and $C$ are the row and column dimensions, respectively. Suppose we have $K$ fine-tuned expert models, each specialized in a distinct downstream task $k \in \{1, \dots, K\}$. The specialized expert weight matrix for layer $l$ and task $k$ is $W_k^{(l)} \in \mathbb{R}^{R \times C}$. 

For each expert and layer, we define the corresponding \textit{task vector} $V_k^{(l)}$ as the parameter difference from the base model:
\begin{equation}
V_k^{(l)} = W_k^{(l)} - W_{\text{base}}^{(l)}
\end{equation}

Our goal is to construct a unified parameter representation that can dynamically transcribe these task vectors on-the-fly, sample-by-sample, while avoiding destructive coordinate interference and remaining immune to batch-level ensembling constraints.

\subsection{1. Carrier Key Generation}

Drawing from hyperdimensional computing~\cite{kanerva2009hyperdimensional}, we generate frozen, random bipolar carrier keys that act as orthogonal spatial reference vectors. For each task $k$ and layer $l$, we first sample two independent random vectors of dimensions $R$ and $C$, respectively:
\begin{align}
r_k^{(l)} &= \text{sign}(\epsilon_R), \quad \epsilon_R \sim \mathcal{N}(0, I_R) \\
c_k^{(l)} &= \text{sign}(\epsilon_C), \quad \epsilon_C \sim \mathcal{N}(0, I_C)
\end{align}
where $r_k^{(l)} \in \{-1, 1\}^R$ and $c_k^{(l)} \in \{-1, 1\}^C$. The 2D spatial carrier key matrix $K_k^{(l)}$ is computed as the outer product of these vectors:
\begin{equation}
K_k^{(l)} = r_k^{(l)} \left( c_k^{(l)} \right)^T \in \{-1, 1\}^{R \times C}
\end{equation}

By construction, since $\mathbb{E}[(r_j^{(l)})^T r_k^{(l)}] = 0$ and $\mathbb{E}[(c_j^{(l)})^T c_k^{(l)}] = 0$ for $j \neq k$, these 2D bipolar carrier keys are mutually pseudo-orthogonal in high-dimensional spaces:
\begin{equation}
\mathbb{E} \left[ \sum_{r, c} K_{j, r, c}^{(l)} K_{k, r, c}^{(l)} \right] = 0 \quad \text{for } j \neq k
\end{equation}

These carrier keys are frozen during training and inference. While we focus here on 2D parameter matrices (linear layers), EHPB generalizes seamlessly to higher-order tensors (such as Conv2D or attention layers) using structured CP or Tucker carrier key decompositions (see Appendix~\ref{app:tensor_decompositions} for a mathematical formulation).

\subsection{2. Holographic Superposition (Parameter Binding)}

Next, we modulate each task-specific parameter offset $V_k^{(l)}$ with its respective spatial carrier key using element-wise Hadamard multiplication. We then superimpose (sum) all bound task vectors into a single physical holographic parameter matrix:
\begin{equation}
W_{\text{holo}}^{(l)} = \sum_{j=1}^K V_j^{(l)} \odot K_j^{(l)}
\end{equation}
where $\odot$ denotes the element-wise Hadamard product. This unified holographic matrix $W_{\text{holo}}^{(l)} \in \mathbb{R}^{R \times C}$ acts as the central superimposed parameter substrate where all expert task vectors are bound and stored in a dormant, non-interfering state.

\subsection{3. Dynamic Transcription Gating (Routing)}

To perform selective dynamic activation, we employ a lightweight, single-layer linear routing network. Let $z(x)_b \in \mathbb{R}^D$ be the globally pooled feature representation extracted from the backbone for sample $b \in \{1, \dots, B\}$ in a batch of size $B$. The routing network computes a sample-wise task-affinity distribution $\alpha_b \in \mathbb{R}^K$:
\begin{align}
s_{k, b} &= W_{\text{router}} z(x)_b + b_{\text{router}} \\
\alpha_{k, b} &= \text{Softmax}\left( s_{b} \right)_k
\end{align}
where $W_{\text{router}} \in \mathbb{R}^{K \times D}$ and $b_{\text{router}} \in \mathbb{R}^K$ are the router's parameters. To ensure high-quality task routing and prevent over-fitting, the router is trained using AdamW with $L_2$ weight decay regularization ($\lambda_{\text{wd}} = 10^{-3}$) on a lightweight 64-sample multi-task calibration set. The complete routing network details, optimization hyper-parameters, offline post-hoc routing calibration, and an empirical sensitivity sweep over calibration data budget sizes are detailed in Appendix~\ref{app:routing_optimizer_details}.

\subsection{4. Holographic Demodulation (Parameter Transcription)}

Rather than averaging ensembling coefficients across the batch dimension, EHPB constructs a sample-specific \textit{unbinding operator} $U_b^{(l)} \in \mathbb{R}^{R \times C}$ for each sample $b$ in the batch:
\begin{equation}
U_b^{(l)} = \sum_{k=1}^K \alpha_{k, b} K_k^{(l)}
\end{equation}

The active, sample-specific weight matrix $W_b^{(l)}$ is transcribed by taking the element-wise Hadamard product of the holographic parameter matrix $W_{\text{holo}}^{(l)}$ and the unbinding operator $U_b^{(l)}$, and adding it to the pre-trained base weight:
\begin{equation}
W_b^{(l)} = W_{\text{base}}^{(l)} + W_{\text{holo}}^{(l)} \odot U_b^{(l)}
\end{equation}

This demodulation is performed independently and in parallel for each sample in the batch, which can be executed highly efficiently using PyTorch's vectorized map ($\mathtt{torch.vmap}$), completely bypassing the hardware constraint of averaging ensembling coefficients across the batch dimension.

\subsection{Mathematical Proof of Reconstruction Accuracy}

We prove that EHPB reconstructs the target dynamic ensembled weight matrix up to a high-frequency, zero-mean cross-talk noise term. 

\begin{theorem}
For any layer $l$ and sample $b$, the holographic demodulated parameter offset $W_{\text{holo}}^{(l)} \odot U_b^{(l)}$ satisfies:
\begin{equation}
W_{\text{holo}}^{(l)} \odot U_b^{(l)} = \sum_{k=1}^K \alpha_{k, b} V_k^{(l)} + \Xi_b^{(l)}
\end{equation}
where $\Xi_b^{(l)}$ is a high-frequency, zero-mean cross-talk noise matrix.
\end{theorem}

\begin{proof}[Proof Sketch]
We expand the unbinding product:
\begin{equation}
W_{\text{holo}}^{(l)} \odot U_b^{(l)} = \left( \sum_{j=1}^K V_j^{(l)} \odot K_j^{(l)} \right) \odot \left( \sum_{k=1}^K \alpha_{k, b} K_k^{(l)} \right)
\end{equation}
Expanding this expression separates the product into a diagonal signal term where $j = k$, and a cross-talk noise term where $j \neq k$. Since the carrier keys are bipolar ($K_k^{(l)} \in \{-1, 1\}^{R \times C}$), their self-product is $K_k^{(l)} \odot K_k^{(l)} = \mathbf{1}_{R \times C}$, which perfectly recovers the target signal $\sum \alpha_{k, b} V_k^{(l)}$. The cross-talk noise matrix $\Xi_b^{(l)} = \sum_{j \neq k} \alpha_{k, b} V_j^{(l)} \odot (K_j^{(l)} \odot K_k^{(l)})$ has an expected value of zero because the carriers are pseudo-orthogonal, i.e., $\mathbb{E}[K_j^{(l)} \odot K_k^{(l)}] = \mathbf{0}_{R \times C}$ for $j \neq k$. See Appendix E for the complete step-by-step mathematical derivation.
\end{proof}

\subsection{The Non-Linearity Confounder in Multi-Layer Noise Propagation}
\label{subsec:non_linearity_confounder}

While Theorem 1 proves that the unbinding cross-talk noise $\Xi_b^{(l)}$ has a zero mean at any individual layer $l$, this property does not guarantee noise-free propagation through a deep neural network. A deep model is highly non-linear, containing non-linear activation functions (e.g., ReLU), Layer Normalization, and Softmax attention. These non-linear transformations do not preserve the zero-mean property of noise, introducing systematic representation biases and signal attenuation that accumulate and explode across deep layers.

We identify and mathematically formalize two highly destructive phenomena (deconstructed in detail in Appendix B):
\begin{enumerate}
\item \textbf{Systematic Positive Bias Rectification:} When zero-mean Gaussian noise is passed through a Rectified Linear Unit ($\text{ReLU}(x) = \max(0, x)$), the negative noise coordinates are clipped to zero while the positive coordinates are preserved. This systematically rectifies the noise into a strictly positive coordinate bias vector $B^{(l)} \geq 0$ (with equality holding only if noise variance is zero). This positive bias offset shifts the mean of propagated representations, compounding across layers, distorting latent representations, and triggering representational collapse.
\item \textbf{LayerNorm Exponential Signal Attenuation:} Weight-reconstruction noise increases pre-activation variance ($\sigma_y^2 \approx \sigma_{\text{target}}^2 + \sigma_e^2$), causing Layer Normalization to scale down the active semantic signal by an attenuation factor of $\eta = \frac{\sigma_{\text{target}}}{\sqrt{\sigma_{\text{target}}^2 + \sigma_e^2}} < 1$. Under our Hadamard binding, $\sigma_e^2 \approx 3 \sigma_{\text{target}}^2$, yielding an attenuation factor of $\eta \approx 0.50$. Across $L=14$ layers, this attenuation compounds exponentially: $\eta^{14} \approx (0.50)^{14} \approx 6 \times 10^{-5}$, completely extinguishing the clean semantic signal.
\end{enumerate}
See Appendix B for the complete mathematical derivations of these non-linear noise propagation effects.

\subsection{Stabilizing Weight Superposition: Residual-EHPB}
\label{subsec:residual_ehpb}

To reconcile the high-frequency noise of hyperdimensional superposition with the coordinate-sensitivity of deep architectures, we propose a hybrid ensembling framework: \textbf{Residual-EHPB}. 

We hypothesize that a deep network's representational capacity is heavily governed by a small, highly critical fraction of its parameter coordinates. To exploit this property, we define a sparse coordinate-wise binary mask $M^{(l)} \in \{0, 1\}^{R \times C}$ that designates the top $p\%$ (e.g., $p = 5\%$) most critical weight coordinates in layer $l$, where coordinate importance is measured by the magnitude of the task vectors:
\begin{equation}
M_{r, c}^{(l)} = \mathbb{I}\left[ \left| \sum_{k=1}^K V_{k, r, c}^{(l)} \right| \geq \tau^{(l)} \right]
\end{equation}
where $\tau^{(l)}$ is a threshold chosen to satisfy $\sum_{r, c} M_{r, c}^{(l)} = p\% \times (R \times C)$. 

In Residual-EHPB, the critical parameters identified by $M^{(l)}$ bypass holographic superposition entirely and are stored with perfect coordinate integrity in an uncompressed sparse matrix, while the remaining $(100 - p)\%$ of parameters are superimposed:
\begin{equation}
W_{\text{holo}, \text{res}}^{(l)} = \sum_{j=1}^K V_j^{(l)} \odot K_j^{(l)} \odot \left( \mathbf{1}_{R \times C} - M^{(l)} \right)
\end{equation}

During test-time demodulation, the active weights are reconstructed as:
\begin{equation}
W_b^{(l)} = W_{\text{base}}^{(l)} + W_{\text{holo}, \text{res}}^{(l)} \odot U_b^{(l)} + \sum_{k=1}^K \alpha_{k, b} V_k^{(l)} \odot M^{(l)}
\end{equation}

By isolating the most critical coordinates, Residual-EHPB guarantees that the corresponding cross-talk noise is zero for these coordinates: $M^{(l)} \odot \Xi_b^{(l)} = \mathbf{0}_{R \times C}$. This provides a mathematically noise-free "clean path" for activation propagation and gradient flow. We must, however, analyze the parameter overhead of Residual-EHPB as the number of tasks $K$ scales. Storing uncompressed parameters for the top $p\%$ coordinates across $K$ expert models requires an active storage footprint scaling as $O(K \times p\% \times P)$. When scaling to a very large portfolio of experts (e.g., $K \geq 100$), this linear scaling with $K$ can become a deployment bottleneck. To mitigate this and keep Residual-EHPB highly parameter-efficient under large $K$, three optimization strategies are proposed: (1) \textbf{Shared Union Gating:} Instead of using task-specific independent masks, we can construct a single globally shared critical coordinate mask $M = \bigvee_{k=1}^K M_k$, or select coordinates that are globally vital across all tasks. (2) \textbf{Manifold Coordinate Overlap:} Because task-specific weights are fine-tuned from a shared initialization, they reside on low-dimensional manifolds and their critical coordinate updates heavily overlap. In practice, the union of critical coordinates grows far slower than $O(K)$. (3) \textbf{Adaptive Sparsity Budgets:} We can dynamically reduce the individual task-wise budget $p_k\%$ as $O(1/K)$ to enforce a hard global uncompressed storage ceiling (e.g., capping total residual parameters at $15\%$ of $P$). Thus, Residual-EHPB remains a highly viable and storage-efficient framework even for massive expert portfolios. We propose Residual-EHPB as a promising, highly practical roadmap to bridge the gap between hyperdimensional weight superposition and coordinate-wise deep neural network performance.

\subsection{The Post-Hoc Model Ensembling Trilemma}
\label{subsec:ensembling_trilemma}

An explicit mathematical expansion of EHPB's dynamic parameter demodulation reveals its exact relationship to standard additive model merging:
\begin{equation}
W_b^{(l)} = W_{\text{base}}^{(l)} + \sum_{k=1}^K \alpha_{k, b} V_k^{(l)} + \Xi_b^{(l)}
\end{equation}

If the cross-talk noise term $\Xi_b^{(l)}$ is set to zero, this expression collapses precisely to a standard, noise-free sample-wise ensembled weight matrix. This equivalence highlights a fundamental question: \textit{Why employ hyperdimensional parameter binding and incur reconstruction noise when one can execute noise-free sample-wise ensembling directly?}

To resolve this conceptual question, we propose a novel theoretical framework: the **Post-Hoc Model Ensembling Trilemma**. We argue that any post-hoc model merging method operating on edge deployment hardware is constrained by three fundamental, mutually competing desiderata:
\begin{enumerate}
\item \textbf{Dynamic Adaptability (D):} The capability to customize parameters on-the-fly, adapting model behavior sample-by-sample based on input-dependent representations.
\item \textbf{Resource Efficiency (R):} The capability to execute multi-task inference with active parameter scaling independent of the number of experts ($O(P)$ active memory footprint, where $P$ is parameter count).
\item \textbf{Weight Integrity (W):} The preservation of native expert parameter coordinates without incurring cross-talk or weight reconstruction noise.
\end{enumerate}

Under this trilemma, no existing post-hoc model merging method can satisfy all three properties simultaneously (Figure~\ref{fig:ensembling_trilemma}).

\begin{figure}[t]
\centering
\begin{tikzpicture}[
    scale=0.92, every node/.style={transform shape},
    vertex/.style={circle, draw, fill=blue!10!white, thick, minimum size=2.5cm, align=center, font=\bfseries\small},
    edge_label/.style={draw, rectangle, rounded corners, fill=green!10!white, thick, font=\small, align=center}
]
    % Coordinates of the triangle vertices
    \coordinate (A) at (0, 3.2);  % Top vertex: Weight Integrity
    \coordinate (B) at (-3.4, -2.0); % Bottom-left vertex: Dynamic Adaptability
    \coordinate (C) at (3.4, -2.0);  % Bottom-right vertex: Resource Efficiency
    
    % Draw the soft background triangle
    \fill[gray!5!white, draw=gray!40!white, line width=1.5pt, dashed] (A) -- (B) -- (C) -- cycle;
    
    % Place the vertices (desiderata)
    \node (WI) at (A) [vertex, fill=teal!10!white, draw=teal] {Weight Integrity\\(W)\\\textit{No Reconstruction}\\\textit{Noise ($\Xi = 0$)}};
    \node (DA) at (B) [vertex, fill=red!10!white, draw=red] {Dynamic\\Adaptability (D)\\\textit{Sample-wise}\\\textit{gating $\alpha(x)$}};
    \node (RE) at (C) [vertex, fill=purple!10!white, draw=purple] {Resource\\Efficiency (R)\\\textit{Active size is $O(P)$}\\\textit{independent of $K$}};
    
    % Place the ensembling methods on the edges
    % Between WI and RE: Static Merging (e.g., Uniform, TIES)
    \node (static) at ($(A)!0.5!(C)$) [edge_label, fill=gray!10!white, draw=gray, xshift=0.3cm, yshift=0.3cm] {Static Merging\\(e.g., Uniform, TIES, soups)\\$O(P)$ size, no noise, static};
    
    % Between WI and DA: Direct Sample-wise Routing
    \node (direct) at ($(A)!0.5!(B)$) [edge_label, fill=blue!10!white, draw=blue, xshift=-0.3cm, yshift=0.3cm] {Direct Sample-wise Routing\\(e.g., vmap-Linear-Router)\\Sample-wise, no noise, $O(K P)$ RAM};
    
    % Between DA and RE: Holographic Parameter Superposition
    \node (ehpb) at ($(B)!0.5!(C)$) [edge_label, fill=orange!15!white, draw=orange!80!black, yshift=-0.5cm] {\textbf{Holographic Superposition (Ours)}\\(EHPB)\\Sample-wise, $O(P)$ memory, but sacrifices W ($\Xi_b \neq 0$)};

\end{tikzpicture}
\caption{The Post-Hoc Model Ensembling Trilemma. Static Merging achieves Resource Efficiency and Weight Integrity but lacks Dynamic Adaptability. Direct Sample-wise Routing (e.g., $\mathtt{vmap}$-Linear-Router) achieves Dynamic Adaptability and Weight Integrity but lacks Resource Efficiency. Our proposed EHPB represents a pioneer in the third corner, achieving Dynamic Adaptability and Resource Efficiency by sacrificing Weight Integrity (incurring cross-talk noise).}
\label{fig:ensembling_trilemma}
\end{figure}

\textbf{1. Static Merging (e.g., Uniform, TIES, Model Soups):}
Static methods average expert weights post-hoc. They achieve **Resource Efficiency** ($O(P)$ memory) and **Weight Integrity** (zero reconstruction noise), but completely lack **Dynamic Adaptability** (their weights are fixed and cannot adapt sample-by-sample).

\textbf{2. Direct Sample-wise Routing (e.g., vmap-Linear-Router):}
Direct sample-wise routing executes vectorized ensembling on-the-fly. It achieves **Dynamic Adaptability** and **Weight Integrity** (zero noise), but completely lacks **Resource Efficiency**. It requires keeping all $K$ expert models active in GPU/edge VRAM simultaneously, leading to an active memory footprint of $O(K \times P)$, which is physically prohibitive for on-device deployment under large $K$ or large parameter size.

\textbf{3. Holographic Parameter Superposition (EHPB):}
Our proposed EHPB represents a pioneer in the third corner of the trilemma. By superimposing all expert models into a single holographic matrix $W_{\text{holo}}$ and dynamically demodulating them via lightweight 1D keys, EHPB achieves **Dynamic Adaptability** and **Resource Efficiency** ($O(P)$ active RAM footprint), but sacrifices **Weight Integrity** by incurring the high-frequency cross-talk noise $\Xi_b^{(l)}$.

\subsection{Resolving the Active Memory Footprint Paradox via Kernel Fusion}
\label{subsec:kernel_fusion}

A subtle technical challenge arises when considering naive parallel batch execution of EHPB in eager-mode deep learning frameworks like PyTorch. In eager-mode, evaluating $W_b = W_{\text{base}} + W_{\text{holo}} \odot U_b$ via vectorized loops ($\mathtt{torch.vmap}$) materializes and stores a batch of sample-specific weight matrices of shape $B \times R \times C$ (or $B \times P$ parameters) in active memory. For a batch size of $B=256$, this eager materialization scales intermediate memory to $O(B \times P)$ during the forward pass, which is larger than the $O(K \times P)$ active storage of keeping all $K=4$ experts in RAM.

However, this is an eager-evaluation artifact rather than an architectural limitation. Compiled on-device deployment pipelines can completely bypass this memory paradox through **Kernel Fusion and Register-Level Demodulation** (e.g., custom Triton or CUDA kernels). As deconstructed in Appendix D, the fused kernel loads only the shared base model $W_{\text{base}}$ and holographic matrix $W_{\text{holo}}$ block-by-block into the GPU's high-speed local cache (SRAM), demodulating weights on-the-fly inside local SM registers:
\begin{equation}
y_{b, r} = \sum_{c=1}^C \left( W_{\text{base}, r, c} + W_{\text{holo}, r, c} \cdot U_{b, r, c} \right) h_{b, c}
\end{equation}
This fused execution prevents writing intermediate weights $W_b$ back to global device memory (HBM), preserving EHPB's true **$O(P)$ active global weight memory footprint** independent of the batch size $B$. See Appendix D for a detailed Triton register allocation and hardware tiling formulation, including local memory and register pressure analyses, and Appendix C for a discussion of the computation-versus-memory-bandwidth trade-offs under this fused register demodulation strategy. Furthermore, we deconstruct the systems-level compatibility of EHPB superposition with low-bit (INT4/INT8) quantized models, proving the Precision Underflow Paradox and formulating hardware-level mitigations in Appendix~\ref{app:quantization_boundaries}.

This formalized trilemma deconstructs the false storage-noise dichotomy, proving that EHPB is a mathematically valid and conceptually necessary architectural exploration. Our work establishes that to bridge the final gap and achieve lossless post-hoc merging on edge devices, future research must focus on identifying hyperdimensional operators that minimize $\Xi_b^{(l)}$ while maintaining the Resource Efficiency and Dynamic Adaptability axes.


